\section{Introduction}

Many physical, biological, and environmental phenomena are subject to uncertainty and random fluctuations. 
Examples include pollutant transport in rivers and groundwater, heat transfer in heterogeneous media, 
atmospheric dispersion, and climate dynamics. In such situations, deterministic partial differential 
equations may fail to accurately represent the variability observed in real systems.

Stochastic partial differential equations (SPDEs) provide a natural framework for incorporating uncertainty 
into mathematical models. By combining differential operators with stochastic forcing terms, SPDEs enable 
the description of systems influenced by random perturbations. Over the last decades, SPDEs have become an 
important research area at the intersection of applied mathematics, numerical analysis, probability theory,
and scientific computing.

The mathematical analysis of stochastic partial differential equations has been extensively developed over the past decades
\cite{DaPrato2014,LiuRockner2015,LordPowellShardlow2014}.
Classical references for finite element methods and numerical approximation of partial differential equations include
\cite{BrennerScott2008,Thomee2006,QuarteroniValli2008,Evans2010,LarssonThomee2003}.
Several numerical methods have been proposed for stochastic partial differential equations, including finite element approximations and time discretization schemes
\cite{Printems2001,KovacsLarssonLindgren2015,CarelliProhl2012}. 
These works provide the theoretical foundations for the study of stochastic evolution equations in 
infinite-dimensional spaces.

From the numerical viewpoint, finite element methods have become one of the most successful techniques for 
the approximation of deterministic and stochastic partial differential equations. Classical references 
include the works of Brenner and Scott \cite{BrennerScott2008} and Thomée \cite{Thomee2006}, which provide 
a rigorous framework for the analysis of Galerkin approximations of elliptic and parabolic problems.

More recently, considerable attention has been devoted to the numerical approximation of stochastic partial 
differential equations. Various spatial discretization techniques, including finite element and finite 
difference methods, have been proposed for stochastic diffusion and reaction--diffusion equations. 

Stability properties, strong convergence, and weak convergence estimates have been established for several
classes of SPDEs.

Despite these advances, stochastic advection--diffusion equations remain comparatively less explored, 
particularly in the context of finite element approximations combining transport effects, diffusion 
mechanisms, and stochastic forcing. Such models are of particular interest in environmental sciences, 
climate modeling, groundwater transport, and uncertainty quantification, where random perturbations may 
significantly influence the system dynamics.

The present work aims to contribute to this research direction by studying a stochastic advection--diffusion 
equation driven by additive noise. In contrast with many existing studies focusing primarily on 
diffusion-dominated stochastic models, the present formulation explicitly incorporates both transport 
and diffusion mechanisms. Furthermore, a complete finite element framework is developed, including the 
derivation of a weak formulation, stability estimates, existence and uniqueness results, and a convergence 
analysis of the fully discrete approximation.


In this work, we consider the stochastic advection--diffusion equation

\begin{equation}
du
=
\left(
-\mathbf b\cdot\nabla u
+
\nabla\cdot(\kappa\nabla u)
+
f
\right)dt
+
\sigma \phi(x)\,dW_t
\label{eq:intro_model}
\end{equation}


in a bounded domain $\Omega \subset \mathbb R^2$, supplemented with homogeneous Dirichlet boundary conditions and a prescribed initial condition. Here, $\mathbf b$ denotes a velocity field, $\kappa$ is the diffusion coefficient, $f$ is a source term, and the stochastic forcing is represented by a Wiener process.

The main objective of this paper is to investigate the mathematical properties of this problem and to develop a finite element approximation. More precisely, we establish a variational formulation of the stochastic advection--diffusion equation, derive energy estimates, prove existence and uniqueness of weak solutions, and study a finite element discretization based on continuous piecewise linear elements and an implicit Euler scheme in time.

The main contributions of this work are summarized as follows:

\begin{itemize}
\item We formulate a stochastic advection--diffusion equation with additive Gaussian noise within a variational framework.
\item We establish stability estimates and prove the existence and uniqueness of weak solutions.
\item We analyze a finite element approximation combined with an implicit Euler time discretization and derive convergence results.
\item We validate the theoretical findings through numerical experiments performed with FEniCSx, including spatial and temporal convergence studies.
\item We investigate the statistical properties of the stochastic solution using Monte Carlo simulations and variance analysis.
\end{itemize}

The present work contributes to this research area by investigating a stochastic advection--diffusion 
equation driven by additive Gaussian noise. The proposed model combines transport, diffusion, and 
stochastic forcing within a unified variational framework. Particular attention is devoted to the 
mathematical analysis of the problem and to the development of a reliable finite element approximation.

More specifically, we derive a weak formulation of the stochastic problem, establish stability estimates, 
prove existence and uniqueness of weak solutions, and analyze a finite element discretization based on 
continuous piecewise linear elements in space and an implicit Euler scheme in time. Numerical experiments 
are performed using FEniCSx in order to validate the theoretical convergence results and to investigate 
the statistical properties of the stochastic solution.

The remainder of the paper is organized as follows. Section~2 introduces the mathematical setting and recalls 
the main functional spaces. Section~3 presents the weak formulation of the stochastic advection--diffusion 
equation. Stability, existence, uniqueness, and convergence results are established in Sections~4--6. 
Section~7 is devoted to the finite element implementation and numerical experiments. Finally, conclusions 
and perspectives are given in Section~8.
